% Chapter 1: Introduction

\chapter{Introduction}
\label{ch:introduction}

\section{Why rationale faithfulness matters}
\label{sec:intro-motivation}

Vision-language models (VLMs) can answer questions about images and explain
their answers in natural language. Instruction-tuned systems such as the Large
Language and Vision Assistant (LLaVA) \citep{DBLP:conf/nips/LiuLWL23a} and
Qwen2.5-VL \citep{DBLP:journals/corr/abs-2502-13923} make these rationales easy
to obtain.
The text can help a reader understand what the model claims to have seen and why
it selected an answer. Its usefulness depends on whether it reflects the
evidence that affected the model's output.

Plausibility and faithfulness are different properties. A plausible rationale
offers a reasonable account to a reader. A faithful rationale reflects the
information that actually contributed to the prediction
\citep{DBLP:conf/acl/JacoviG20}. Studies of language models have shown that a
cue can change an answer without appearing in the generated chain of thought
\citep{DBLP:conf/nips/TurpinMPB23}. Similar concerns arise in visual question
answering (VQA), where question patterns and shallow regularities between questions and
objects can support predictions even when the image contributes less than
expected \citep{DBLP:conf/iccv/DancetteCTC21,DBLP:conf/emnlp/SiMZLL0CWZ22}.

This thesis uses a two-stage answer and rationale pipeline. Stage~1 predicts one
of four answers. Stage~2 receives a supplied answer and generates a rationale
for it. Separating the stages lets us ask whether the answer changes when visual
evidence changes and whether the rationale responds to the same intervention.
For the primary Stage~2 experiment, we hold the supplied gold answer fixed so
that an image intervention cannot also change the answer that Stage~2 is asked
to justify. The gold answer is the correct choice recorded by the dataset's
human annotators.

Chain-of-thought (CoT) prompting elicits intermediate reasoning text before an
answer \citep{WeiEtAl22CoT}. In this thesis, the generated output is divided
into a \texttt{<reasoning>} span and a \texttt{<final>} span, and that division
is the thesis's own output convention. The
\texttt{<reasoning>} part is step-by-step explanatory text. The
\texttt{<final>} part is a shorter justification written for a reader. The
primary RQ2 measure uses the \texttt{<final>} span. That span is the concise
output the model presents as its justification. It is recovered by one fixed parsing rule, and it is the unit
specified for the primary analysis before the results were seen. We also
analysed the reasoning span and the complete generation in
Appendix~\ref{sec:app-rq2-sensitivity} to check whether the broad ordering of
the interventions depends on using the final span. Generated CoT text is a model
output. It does not directly record the computation that produced an answer
\citep{DBLP:journals/corr/abs-2307-13702}.

Faithfulness concerns whether an explanation reflects the information actually
used by the model. This study measures behavioural changes in Stage~1 answers
and Stage~2 rationale text under controlled alterations of visual evidence.
Chapter~\ref{ch:problem-statement} sets out what that evidence licenses.

\section{Study design}
\label{sec:intro-design}

We used Visual Commonsense Reasoning (VCR), a multiple-choice benchmark built
from movie scenes \citep{DBLP:conf/cvpr/ZellersBFC19}. The evaluation population
contains 2,653 question records drawn from 2,400 source frames. We adapted the
same Qwen2.5-VL-3B-Instruct base model separately for four input arms:

\begin{description}
\item[\texttt{plain}] the source image without an added description;
\item[\texttt{point}] the image with labelled polygon overlays;
\item[\texttt{plain\_desc}] the source image with a fixed description; and
\item[\texttt{point\_desc}] the polygon image with the same description.
\end{description}

The arms vary two possible aids. A polygon overlay points to referenced
people or objects in the image. A description provides scene information in
text. We evaluated every arm under six image conditions: source, grey,
mismatch, random pixel mask, Gaussian noise, and horizontal mirror. The same
mismatch map, which pairs records from different source frames, was used at
both stages.

Stage~1 is evaluated through accuracy, gaps relative to source, and movements
between correct and wrong predictions. Stage~2 is evaluated through BERTScore
Drift relative to source, which measures change in generated rationale text.
Three comparisons help interpret the Drift scale:

\begin{itemize}
\item \textbf{Mirror.} The image is reversed horizontally while its scene
content is preserved.
\item \textbf{\mbox{Wrong-answer} control.} The source image remains fixed, but
Stage~2 receives an answer that differs from the VCR gold answer. Here, wrong
means different from the dataset's gold answer.
\item \textbf{Reference using rationales from different records.} Within each
arm, a rationale
generated under the source condition for one record is paired with a rationale
generated under the source condition for a different record, drawn from a
different source frame. The resulting Drift provides an empirical reference for
unrelated rationale texts.
\end{itemize}

\section{Research questions}
\label{sec:intro-rqs}

The thesis addresses exactly two research questions:

\begin{description}
\item[RQ1.] How does Stage~1 answer behaviour change when visual evidence is altered or removed across the four input arms?
\item[RQ2.] When the supplied gold answer is held fixed, how do Stage~2 rationales change when visual evidence is altered or removed across the four input arms?
\end{description}

Together, the two questions provide behavioural evidence relevant to rationale
faithfulness by examining answer dependence and rationale sensitivity under the
same controlled changes in visual evidence. Chapter~\ref{ch:discussion}
interprets the two sets of results together.

\section{Contributions}
\label{sec:intro-contributions}

This thesis makes four contributions.

\begin{itemize}
\item It gives a matched evaluation of Stage~1 answers under six image conditions across four
input arms on the same 2,653 records.
\item It provides the complete primary Stage~2 rationale set conditioned on the
gold answer across the source image and five interventions.
\item It interprets BERTScore Drift against three empirical references while
reporting parser coverage and retaining the raw metric specified in the study plan.
\item It compares answer and rationale patterns associated with descriptions
and polygon overlays across both stages.
\end{itemize}

\section{Thesis structure}
\label{sec:intro-overview}

Chapter~\ref{ch:problem-statement} defines the claim boundary and scope of the
study.
Chapter~\ref{ch:related-work} connects the study to VQA shortcuts, rationale
generation, and behavioural faithfulness tests. Chapters~\ref{ch:setup}
and~\ref{ch:framework} describe the experiment and analysis. Chapter~\ref{ch:results}
answers the two research questions. Chapter~\ref{ch:discussion} interprets the
two stages together, and Chapter~\ref{ch:conclusion} concludes the thesis.
